\subsection{局部紧域上的模函数}

设 K 为局部紧域（不必交换）。回顾，函数 mod（或 mod \(_K\)）已在 K 上定义如下（积分，第 VII 章，第 1 节，第 10 节，定义 6）：mod \(_K(0) = 0\)，对于 K 中 \(x \neq 0\)，数 mod \(_K(x)\) 是 K 的加法群的自同构 \(y \mapsto xy\) 的模。

命题 1. 若 K 是局部紧域，则函数 mod \(_{K}\) 属于 \(\mathcal{V}(K)\)（\(\S\) 6，第 1 节）。此外：

\begin{enumerate}
\def\labelenumi{(\roman{enumi})}
\item
  若 \(s > 0\) 使得 \((\mathrm{mod}_{\mathbf{K}})^s = g\) 是绝对值，则 \(g\) 定义 \(\mathbf{K}\) 上的拓扑。
\item
  若 K 非离散且 \(mod_{K}\) 是超度量绝对值，则 K 上存在一个赋范离散赋值 v，其环是紧的且剩余域是含 q 个元素的有限域，从而 \(mod_{K} = q^{-v}\)。K 上的拓扑由 v 定义。
\end{enumerate}

这由 § 6，第 1 节，命题 1；§ 5，第 1 节，命题 2；以及积分，第 VII 章，第 1 节，第 10 节，命题 12 和 13 得到。

命题 2. 设 K、K' 是两个（不必交换的）局部紧域，K 是 K' 的拓扑子域且 K 非离散。则：

\begin{enumerate}
\def\labelenumi{(\roman{enumi})}
\item
  \(\mathbf{K}'\) 是 \(\mathbf{K}\) 上的有限维左（分别地，右）向量空间。
\item
  如果 \(\mathbf{K}\) 包含于 \(\mathbf{K}'\) 的中心中，则对所有 \(x \in \mathbf{K}'\)，
\end{enumerate}

\[
\operatorname{mod} _ {\mathbf {K} ^ {\prime}} (x) = \operatorname{mod} _ {\mathbf {K}} (\mathrm{N} _ {\mathbf {K} ^ {\prime} / \mathbf {K}} (x)).\tag{1}
\]

由于 K 是非离散的完备赋值域，断言 (i) 由《拓扑向量空间》第一章，第 2 节，第 4 节，定理 3 得到；断言 (ii) 正是《积分》第七章，第 1 节，第 11 节，命题 17。

推论 1. 每个中心非离散的局部紧域都是其中心上的有限秩扩张。

局部紧域 K 的中心 Z 在 K 中是闭的，因而是局部紧的。

推论 2. 设 K' 是局部紧域，K 是 K' 的闭子域（两者均不必交换）。若 K' 是 K 上的 n 维左（分别地，右）向量空间，则

\[
\operatorname{mod} _ {\mathrm{K} ^ {\prime}} (x) = (\operatorname{mod} _ {\mathrm{K}} (x)) ^ {n} \quad \text {for all} \quad x \in \mathrm{K}.\tag{2}
\]

一般已知，在 K 上有限维的（左或右）向量空间中，比例为 \(x \in K\) 的伸缩的模等于 \((\text{mod}_{\mathbf{K}}(x))^{n}\)；只需将此应用于 \(K'\) 即可。

